\section{Conclusion}
\label{sec:conclusion}

This paper asked what retrieval strategy, if any, best supports LLM-based
detection of malicious PyPI packages. Using the EASE 2025 dataset, a fixed
\texttt{Llama-3.1-8B-Instruct} generator, and a code-only knowledge base, we
compared a zero-shot baseline against dense, sparse, hybrid, behavior-based,
and contrastive retrieval.

Three conclusions appear robust. Any code-level retrieval is a large win,
lifting malicious-class F1 from 0.888 to 0.97, mostly through recall. The
mechanism matters surprisingly little: plain BM25 matches a 4B embedding
retriever within 0.4 F1 points, so sparse retrieval is a strong, low-cost
default. And contrastive retrieval offers a favorable precision/recall
trade-off: retrieving benign counterexamples cuts the false-positive rate to
0.1\% at 0.948 recall. Conversely, reducing queries to AST behavior summaries
is counterproductive (F1 0.878), demonstrating that signal loss in query
construction, not retrieval per se, is where methods fail.

We also documented the residual failures that no retrieval over
\texttt{setup.py} can fix---metadata-cloned typosquats---and the recurring
false positives over security-adjacent legitimate packages, as guidance for
the next generation of designs, which should extend evidence beyond a single
file, add name-level signals, and calibrate retrieval toward triage
workflows.
\endinput
